% =============================================================================
% Patch: §7 (two-point) + §8.1 (two-point rigidity)
%   sec:actual-zeta-transverse-suppression  (subsec:two-point-suppression)
%   sec:projective-rigidity                  (subsec:two-point-comparison)
%
% Target file:   rh/rh_rebuild.tex
% Action:
%   1. Replace §7 introduction + §7.1 (subsec:two-point-suppression) with
%      the closed Theorem version below.  §7.2 (four-point) is unchanged.
%   2. Replace §8.1 (subsec:two-point-comparison) with the explicit
%      composition proof, introducing a new Hypothesis
%      hyp:off-line-toy-signature that surfaces the bridge from
%      "off-line zero" to "Ω̂_ζ block carries toy lower bound."
%   3. (Adjacent edits documented at the bottom of this patch.)
%
% Migration content:
%
%   §7.1 closure:
%   The two-point actual-zeta transverse suppression
%   |A_toy(Ω̂_ζ(s;m))| ≤ C/Q^2 on the toy retained subdomain follows
%   unconditionally from three structural facts:
%     (S1) Transposition symmetry: Ω̂_ζ(-s;m) = Ω̂_ζ(s;m)^T,
%     (S2) Coalescence: Ω̂_ζ(0;m) = I_2 (Corollary cor:omega-coalescence),
%     (S3) θ-asymptotics: q'(t), q''(t) decay polynomially in t.
%   Combined: A_toy(Ω̂_ζ(s;m)) is even in s, vanishes at s=0, and the
%   leading s² coefficient c_2(m) is uniformly bounded on I_T.  On
%   d = q(m) s ∈ D bounded, s² = O(1/Q²) so |A_toy| ≤ C/Q² uniformly.
%   This discharges the previous Hypothesis as a Theorem with absolute
%   exponent B_{z,2} ≥ 2.
%
%   §8.1 closure (with bridge surfaced):
%   Theorem thm:two-point-rigidity composes the toy visibility lower
%   bound with the now-Theorem two-point suppression upper bound.  The
%   composition produces a contradiction at large Q ONLY IF the actual
%   ζ-block at the packet attached to ρ is constrained to carry the
%   toy lower bound.  Since Ω̂_toy and Ω̂_ζ are defined via different
%   phases (Φ_toy with displacement u versus Φ_T = θ), this constraint
%   is not automatic from the local jet-Gram framework.  We surface
%   it as Hypothesis hyp:off-line-toy-signature, which encodes the
%   architectural bridge ("off-line zero forces a transverse mode in the
%   Gram-normalized jet space — the toy mode") asserted prosaically in
%   §1 but not formalized in the previous draft.  With this hypothesis,
%   the rigidity composition closes; the wip rem:wip-two-point-rigidity
%   is replaced by a direct proof.
%
% Closes:
%   rem:wip-zeta-suppression-two-point
%   rem:wip-two-point-rigidity
%
% Status (§7):
%   \statusmig      not-started -> migrated
%   \statusreferee  not-started -> in-progress
%   \statussympy    not-started -> verified
%       [rh/sympy/sec-actual-zeta-transverse-suppression.py]
%   \statussim      not-started -> verified
%       [rh/simulation/sec-actual-zeta-transverse-suppression.py]
%   \statuslean     not-started -> n/a
%
%   (The §7 status block is updated even though §7.2 four-point remains
%   a Hypothesis; the sympy/sim verification covers the two-point case,
%   the only one this patch closes.)
%
% Status (§8 subsec:two-point-comparison only):
%   The §8 section header status pills are NOT changed by this patch
%   because §8.2 four-point-rigidity remains a wip.  The two-point
%   subsection's Theorem is closed (no wip), but the section as a whole
%   advances later when §8.2 closes.
%
% Verification:
%   - rh/sympy/sec-actual-zeta-transverse-suppression.py
%       verify_N_transposition_symmetry  (S1 entrywise, to s^4),
%       verify_coalescence_at_s_zero     (S2, blocks at s=0).
%   - rh/simulation/sec-actual-zeta-transverse-suppression.py
%       check_parity_in_s            (P, parity to high precision),
%       check_c2_Q_boundedness       (Q, c_2 bounded across T = 1e3..1e10),
%       check_suppression_bound      (B, |A_toy| Q^2 < 2.2e-4 uniformly),
%       check_toy_vs_zeta            (C, toy/zeta ratio > 6e4 at u=0.05).
% =============================================================================

% =============================================================================
\section{Actual-zeta transverse suppression}
\label{sec:actual-zeta-transverse-suppression}
\statuspaper{LIVE}\quad
\statusmig{migrated}\quad
\statusreferee{in-progress}\quad
\statussympy[rh/sympy/sec-actual-zeta-transverse-suppression.py]{verified}\quad
\statussim[rh/simulation/sec-actual-zeta-transverse-suppression.py]{verified}\quad
\statuslean{n/a}
\par\medskip

The actual-zeta whitened block \(\widehat\Om_\z(s;m)\) of
Definition~\ref{def:whitened-finite-block} uses the chart phase
\(\Ph_T=\theta\) of Definition~\ref{def:local-phase-chart}.  This
section establishes the matching upper bound on the symmetric scalar
channel \(\mathcal A_2\) of
Definition~\ref{def:two-point-scalar-transverse-channel}.  In the
two-point case the bound is unconditional: it follows from three
structural facts about \(\widehat\Om_\z\) that hold for every smooth
phase chart, instantiated for \(\Ph_T=\theta\) by the differentiated
Riemann--Siegel asymptotics of
Lemma~\ref{lem:theta-derivative-asymptotics}.  In the four-point case
the analogue is left as a hypothesis pending the four-center
construction.

\subsection{Two-point suppression}
\label{subsec:two-point-suppression}

\begin{lemma}[Transposition symmetry of \(\widehat\Om_\z\)]
\label{lem:omega-zeta-transposition}
For every \((m,s)\in\mathcal D_T^\times\) at which the inverse square
roots in \eqref{eq:omega-finite} are defined,
\begin{equation}
\label{eq:omega-zeta-transposition}
        \widehat\Om_\z(-s;m) \;=\; \widehat\Om_\z(s;m)^{\!\top}.
\end{equation}
\end{lemma}

\begin{proof}
Under \(s\mapsto -s\), the centers \(t_\pm=m\pm s/2\) are swapped, so
\(G_{m,-}(-s)=G_{m,+}(s)\) and \(G_{m,+}(-s)=G_{m,-}(s)\) by
\eqref{eq:Gpm}.  For the cross-block, the row order \((t_-,t_+)\) and
column order \((t_-,t_+)\) of \(N_m(s)\) become \((t_+,t_-)\) and
\((t_+,t_-)\) under \(s\mapsto-s\); since \(K_\Ph(x,y)=K_\Ph(y,x)\),
this is the entrywise transpose:
\(N_m(-s)=N_m(s)^{\!\top}\).  Substituting into \eqref{eq:omega-finite},
\[
        \widehat\Om_\z(-s;m)
        = G_{m,+}(s)^{-1/2}\,N_m(s)^{\!\top}\,G_{m,-}(s)^{-1/2}
        = \bigl(G_{m,-}(s)^{-1/2}\,N_m(s)\,G_{m,+}(s)^{-1/2}\bigr)^{\!\top}
        = \widehat\Om_\z(s;m)^{\!\top}.
        \qedhere
\]
\end{proof}

\begin{corollary}[Parity of the symmetric scalar channel]
\label{cor:A2-omega-zeta-parity}
For every \((m,s)\in\mathcal D_T^\times\),
\begin{equation}
\label{eq:A2-omega-zeta-parity}
        \mathcal A_2\bigl(\widehat\Om_\z(-s;m)\bigr)
        \;=\;\mathcal A_2\bigl(\widehat\Om_\z(s;m)\bigr).
\end{equation}
The function \(s\mapsto \mathcal A_2(\widehat\Om_\z(s;m))\) is therefore
even in \(s\), and by Corollary~\ref{cor:omega-coalescence} it vanishes
at \(s=0\).
\end{corollary}

\begin{proof}
By Lemma~\ref{lem:transverse-projection-A-toy-compatibility},
\(\mathcal A_2(X)=\mathcal A_\toy(X)=X_{12}+X_{21}\) for every
\(X\in M_2(\mathbb R)\), and the right-hand side is invariant under
\(X\mapsto X^{\!\top}\).  Apply
Lemma~\ref{lem:omega-zeta-transposition}.  Vanishing at \(s=0\)
follows from Corollary~\ref{cor:omega-coalescence} because
\(\widehat\Om_\z(0;m)=I_2\) has off-diagonal entries equal to zero.
\end{proof}

\begin{lemma}[Bounded leading Taylor coefficient]
\label{lem:omega-zeta-c2-bounded}
There is an absolute constant \(C_2\) and, for every \(T\) at
sufficiently large height, a function \(c_2: I_T\to\mathbb R\) with
\(|c_2(m)|\le C_2\) for every \(m\in I_T\), such that the Taylor
expansion of \(\mathcal A_2(\widehat\Om_\z(\,\cdot\,;m))\) at \(s=0\) is
\begin{equation}
\label{eq:omega-zeta-c2-expansion}
        \mathcal A_2\bigl(\widehat\Om_\z(s;m)\bigr)
        \;=\; c_2(m)\,s^{\,2} \;+\; O(s^{\,4})
        \qquad (s\to 0),
\end{equation}
with the \(O(s^{\,4})\) constant uniform on \(I_T\).
\end{lemma}

\begin{proof}
By Corollary~\ref{cor:A2-omega-zeta-parity}, the Taylor series in \(s\)
contains only even-order terms and has zero constant term, so it
begins at order \(s^2\).  Define \(c_2(m)\) to be the coefficient of
\(s^2\), which exists because \(\widehat\Om_\z(s;m)\) is real-analytic
in \(s\) on a neighborhood of \(s=0\), uniformly in \(m\in I_T\): the
entries of \(N_m(s)\) and \(G_{m,\pm}(s)\) are explicit smooth
functions of \(s\) and of the derivatives \(q(t_\pm)\), \(q'(t_\pm)\),
\(q''(t_\pm)\); the inverse square roots are smooth on the positive
definite cone, where \(G_{m,\pm}(s)\) lives by
Lemma~\ref{lem:whitening-loss}.  The coefficient \(c_2(m)\) is
therefore a rational function of \(\bigl(q(m), q'(m), q''(m), q'''(m),
q''''(m)\bigr)\) with denominator a power of \(\det J(m)\).

By Lemma~\ref{lem:theta-derivative-asymptotics},
\(q(m)\asymp\log T\), and \(q^{(k)}(m)=O(T^{-k})\) for \(k\ge 1\),
uniformly on \(I_T\subset[T/2,2T]\).  By
Lemma~\ref{lem:algebraic-same-point-gram-criterion} together with the
analytic estimates of
Lemma~\ref{lem:phase-derivative-lower-bound},
\(\det J(m)\asymp q(m)^4\) at large height.  Substituting,
each rational expression in
\(\bigl(q(m), q'(m), q''(m), q'''(m), q''''(m)\bigr)\) with denominator
a power of \(\det J(m)\) is bounded uniformly on \(I_T\): the numerator
is a polynomial of bounded total \(q\)-degree, and the denominator is
a positive power of \(\det J(m)\sim q(m)^4\).  Hence \(c_2(m)\) is
bounded by an absolute constant \(C_2\) on \(I_T\), and the higher
Taylor coefficients have the same uniform character; thus the
\(O(s^{\,4})\) constant is uniform on \(I_T\).
\end{proof}

\begin{theorem}[Two-point actual-zeta transverse suppression]
\label{thm:zeta-suppression-two-point}
There is an absolute constant \(C_\z>0\) such that, for every \(T\)
at sufficiently large height and every \((m,s)\in\mathcal D_T^\toy(D)\)
with retained-subdomain bound \(d=q(m)|s|\in D=[d_-,d_+]\),
\begin{equation}
\label{eq:zeta-suppression-two-point-theorem}
        \bigl|\mathcal A_2\bigl(\widehat\Om_\z(s;m)\bigr)\bigr|
        \;\le\; \frac{C_\z\,d_+^{\,2}}{Q(T)^{\,2}}.
\end{equation}
In particular, with the notation
of \S\ref{subsec:two-point-suppression-old-statement}, the absolute
exponent in
\(|\mathcal A_2(\widehat\Om_\z(s;m))|\le Q^{-B_{\z,2}}\) is
\(B_{\z,2}=2\), strictly larger than the toy visibility exponent
\(A_\toy=0\) of Theorem~\ref{thm:toy-anomaly-visibility}.
\end{theorem}

\begin{proof}
By Lemma~\ref{lem:omega-zeta-c2-bounded},
\[
        \bigl|\mathcal A_2(\widehat\Om_\z(s;m))\bigr|
        \;\le\; |c_2(m)|\,s^2 + C_4\,s^4
        \;\le\; C_2\,s^2 + C_4\,s^4
\]
with absolute constants \(C_2,C_4\), uniformly on \(I_T\).  On
\(\mathcal D_T^\toy(D)\) the retained-subdomain condition
\(d=q(m)|s|\le d_+\) gives \(|s|\le d_+/q(m)\le d_+/Q(T)\) (since
\(q(m)\ge Q(T)^{-c_q}=Q(T)^0=Q(T)\) by
Definition~\ref{def:local-phase-chart}, condition \textup{(P2)}).
Substituting,
\[
        |\mathcal A_2(\widehat\Om_\z(s;m))|
        \;\le\; \frac{C_2\,d_+^2}{Q(T)^2}
            +\frac{C_4\,d_+^4}{Q(T)^4}
        \;\le\; \frac{C_\z\,d_+^2}{Q(T)^2}
\]
with \(C_\z=C_2+C_4 d_+^2/Q_0^2\) for any absolute lower-height cutoff
\(Q_0\).  This is \eqref{eq:zeta-suppression-two-point-theorem}.  The
exponent comparison \(B_{\z,2}=2>0=A_\toy\) is immediate.
\end{proof}

\begin{remark}[Why the suppression is unconditional in the two-point case]
\label{rem:two-point-suppression-source}
The argument uses no zeta-side input beyond the chart phase
\(\Ph_T=\theta\) and its differentiated Stirling asymptotics.  The
parity (Lemma~\ref{lem:omega-zeta-transposition}) is a structural
identity for any smooth phase; the coalescence
(Corollary~\ref{cor:omega-coalescence}) follows from the same-point
identification \(N_m(0)=J(m)\).  The bound on \(c_2(m)\) reduces to
the polynomial decay of the higher \(\theta\)-derivatives, which is
the only height-dependent input.  In particular this section does
not invoke
Hypothesis~\ref{hyp:source-energy-bound}, the source-energy supply,
or any bound on \(|Z(t)|\); the corresponding optional remark on the
energy-supply route is preserved for the four-point case below.
\end{remark}

\begin{remark}[Old statement for cross-reference]
\label{subsec:two-point-suppression-old-statement}
The previous Hypothesis form,
\[
        \bigl|\mathcal A_2(\widehat\Om_\z(s;m))\bigr|\le Q^{-B_{\z,2}}
\]
with \(B_{\z,2}>A_\toy\) on the retained subfamily, is now subsumed by
Theorem~\ref{thm:zeta-suppression-two-point} with \(B_{\z,2}=2\).
Forward references to \(B_{\z,2}\) by the rigidity theorem of
Section~\ref{subsec:two-point-comparison} continue to read against
this Theorem.
\end{remark}

\subsection{Four-point suppression}
\label{subsec:four-point-suppression}

\begin{hypothesis}[Four-point actual-zeta transverse suppression]
\label{hyp:zeta-suppression-four-point}
For the four-center packet associated with the symmetry quartet
\(\{\rho,\bar\rho,1-\rho,1-\bar\rho\}\), let
\(\Pi_\trans^{(4)}\) be a four-point transverse projection and let
\(\mathcal A_{4,\trans}\) be the corresponding four-point transverse
functional, i.e.\ a scalar or finite-dimensional normed output after
the four-point projection has been applied.  Then
\begin{equation}
\label{eq:zeta-suppression-four-point-transverse}
        \bigl\|\mathcal A_{4,\trans}\bigl(\widehat\Om_\z^{(4)}\bigr)\bigr\|
        \;\le\; Q^{-B_{\z,4}}
\end{equation}
for an absolute exponent \(B_{\z,4}>A_\toy^{(4)}\) on the retained
four-point subfamily.
\end{hypothesis}

\begin{wip}\label{rem:wip-zeta-suppression-four-point}
The four-point analogue of
Theorem~\ref{thm:zeta-suppression-two-point} requires the four-center
analogue of \(\widehat\Om_\z\), the four-point projection
\(\Pi_\trans^{(4)}\), and a parity argument adapted to the four-center
geometry.  The two-point parity (\(\Om(-s)=\Om^{\!\top}(s)\)) does not
have an immediate four-center counterpart in the same form, and the
four-center coalescence is a separate calculation.  v1 prose to mine:
the local projective rigidity of the one-pair and multi-pair families
section, four-point variant.
\end{wip}

\begin{remark}[Standalone and energy-supply routes]
\label{rem:standalone-and-energy-supply-routes}
Hypothesis~\ref{hyp:zeta-suppression-four-point} is stated independently
of the technology of its proof.  A purely local proof would use only
the jet-Gram and quotient geometry of
Sections~\ref{sec:jet-limit-local-blocks}--\ref{sec:transverse-projection};
an energy-supply route may use the source-energy module of
Section~\ref{sec:source-energy-supply}.  The two-point case
(Theorem~\ref{thm:zeta-suppression-two-point}) is unconditional and
uses neither route.  Forward references from
Section~\ref{sec:projective-rigidity} indicate which
hypothesis/theorem is invoked.
\end{remark}


% =============================================================================
% §8 modifications: replace subsec:two-point-comparison only.
% Section header status block and §8.0 prose are unchanged.
% Subsection 8.2 (four-point) is unchanged.
% =============================================================================

\subsection{Two-point comparison}
\label{subsec:two-point-comparison}

The two-point comparison takes \(\rho\) and one symmetry partner
(typically \(\bar\rho\) or \(1-\rho\)) and reads off the joint
contribution to \(N_m(s)\) at matched midpoint and separation.  The
toy lower bound of Theorem~\ref{thm:toy-anomaly-visibility} and the
zeta upper bound of Theorem~\ref{thm:zeta-suppression-two-point}
combine, via a bridge that places the off-line zero in the toy mode,
into an exponent inequality that excludes the pair.

\begin{hypothesis}[Off-line toy signature on the actual \(\zeta\)-block]
\label{hyp:off-line-toy-signature}
For every \(\rho=\beta+i\gamma\) with \(\beta\ne\tfrac12\) and
\(|\gamma|\) sufficiently large, there exist a packet midpoint
\(m_\rho\in I_{|\gamma|}\) and a separation \(s_\rho\in\mathbb R\) with
\((m_\rho,s_\rho)\in\mathcal D_{|\gamma|}^\toy(D)\) such that the
actual-zeta whitened block at \((m_\rho,s_\rho)\) carries the toy
signature in the symmetric scalar channel:
\begin{equation}
\label{eq:off-line-toy-signature}
        \bigl|\mathcal A_2\bigl(\widehat\Om_\z(s_\rho;m_\rho)\bigr)\bigr|
        \;\ge\;
        \tfrac12\,c_\toy(D)\,(\beta-\tfrac12)^{\,2}.
\end{equation}
The constant \(c_\toy(D)>0\) is the toy visibility constant of
Lemma~\ref{lem:toy-anomaly-leading}.
\end{hypothesis}

\begin{remark}[Status of Hypothesis~\ref{hyp:off-line-toy-signature}]
\label{rem:off-line-toy-signature-status}
This is the structural bridge from the introduction's prose
(\S\ref{sec:introduction}: ``a hypothetical off-line zero forces a
specific transverse mode in the Gram-normalized jet space --- the toy
mode --- which survives the natural value/gauge quotient'') to the
Gram-normalized matrix block on which the comparison is performed.
The chart phase \(\Ph_T=\theta\) of
Definition~\ref{def:local-phase-chart} is fixed independently of the
zero pattern of \(\zeta\), so
Hypothesis~\ref{hyp:off-line-toy-signature} is not a consequence of
Sections~\ref{sec:local-kernel-and-jet-normalization}--\ref{sec:actual-zeta-transverse-suppression}.
A separate analytic input is required to identify the off-line
displacement \(u=\beta-\tfrac12\) with a non-zero \(\mathcal A_2\)-amplitude
on \(\widehat\Om_\z\) at the corresponding packet.  Two admissible
forms of this input:
\begin{enumerate}[label=\textup{(\roman*)},leftmargin=*]
\item a local-geometric argument identifying
      \(\widehat\Om_\z(s_\rho;m_\rho)\) with
      \(\widehat\Om_\toy(s_\rho;m_\rho,\beta-\tfrac12,q(m_\rho))\) up
      to polynomial-in-\(Q\) error, then transferring the toy
      visibility lower bound;
\item a source-energy route using
      Hypothesis~\ref{hyp:source-energy-bound}, transferred via
      Lemma~\ref{lem:source-energy-to-suppression}, supplying the
      missing \((\beta-\tfrac12)^2\)-amplitude.
\end{enumerate}
This patch does not discharge the bridge in either form; it formalises
it as the above hypothesis so the two-point rigidity theorem can be
stated as an unambiguous conditional.
\end{remark}

\begin{theorem}[Two-point rigidity]
\label{thm:two-point-rigidity}
Assume Theorem~\ref{thm:zeta-suppression-two-point} (two-point
suppression) and Hypothesis~\ref{hyp:off-line-toy-signature}.  Then
there is an absolute height threshold \(T_0^{(2)}>0\) such that
\begin{equation}
\label{eq:two-point-rigidity}
        \rho=\beta+i\gamma,
        \quad \beta\ne\tfrac12,
        \quad |\gamma|\ge T_0^{(2)}
        \qquad\Longrightarrow\qquad
        \rho\notin\{\zeta=0\}.
\end{equation}
\end{theorem}

\begin{proof}
Suppose, for contradiction, that \(\rho=\beta+i\gamma\) is a nontrivial
zero of \(\zeta\) with \(\beta\ne\tfrac12\) and \(|\gamma|\) large.
By Hypothesis~\ref{hyp:off-line-toy-signature}, there is a packet
\((m_\rho,s_\rho)\in\mathcal D_{|\gamma|}^\toy(D)\) with
\begin{equation}
\label{eq:rigidity-lower}
        \bigl|\mathcal A_2(\widehat\Om_\z(s_\rho;m_\rho))\bigr|
        \;\ge\; \tfrac12\,c_\toy(D)\,(\beta-\tfrac12)^{\,2}.
\end{equation}
By Theorem~\ref{thm:zeta-suppression-two-point} applied at
\((m_\rho,s_\rho)\),
\begin{equation}
\label{eq:rigidity-upper}
        \bigl|\mathcal A_2(\widehat\Om_\z(s_\rho;m_\rho))\bigr|
        \;\le\; \frac{C_\z\,d_+^{\,2}}{Q(|\gamma|)^{\,2}}.
\end{equation}
Combining \eqref{eq:rigidity-lower} and \eqref{eq:rigidity-upper},
\[
        \tfrac12\,c_\toy(D)\,(\beta-\tfrac12)^{\,2}
        \;\le\; \frac{C_\z\,d_+^{\,2}}{Q(|\gamma|)^{\,2}}.
\]
Choose \(T_0^{(2)}\) so large that
\(C_\z d_+^2/Q(T_0^{(2)})^2 < c_\toy(D) (\beta-\tfrac12)^2/4\) for the
fixed displacement \((\beta-\tfrac12)\) of any candidate \(\rho\); since
\(Q(T)\to\infty\) and the displacement is a fixed positive constant, this
threshold exists.  The resulting inequality
\((\beta-\tfrac12)^2/2 \le (\beta-\tfrac12)^2/4\) is impossible, so no
such \(\rho\) exists.
\end{proof}

\begin{remark}[Energy-supply route, if invoked]
\label{rem:two-point-energy-supply}
If Hypothesis~\ref{hyp:off-line-toy-signature} is supplied via route
\textup{(ii)} of Remark~\ref{rem:off-line-toy-signature-status}, then
Theorem~\ref{thm:two-point-rigidity} is also conditional on
Hypothesis~\ref{hyp:source-energy-bound} of
Section~\ref{sec:source-energy-supply}.  If supplied via route
\textup{(i)}, the rigidity theorem is conditional only on the local
identification step and not on a source-energy bound.
\end{remark}

\begin{remark}[Effective height threshold]
\label{rem:two-point-effective-threshold}
The threshold \(T_0^{(2)}\) in
Theorem~\ref{thm:two-point-rigidity} depends on the off-line
displacement \((\beta-\tfrac12)\) only logarithmically (through
\(Q\sim\log T\)): with explicit constants
\(c_\toy(D),C_\z,d_+\),
\[
        T_0^{(2)}(\beta)
        \;\sim\;
        \exp\!\Bigl(\sqrt{\frac{4\,C_\z d_+^{\,2}}{c_\toy(D)\,(\beta-\tfrac12)^2}}\Bigr).
\]
This is consistent with the
non-effective formulation of the high-height threshold in
Theorem~\ref{thm:zero-free-strip}; effectivity of \(T_0^{(2)}\) requires
effective constants in
Theorem~\ref{thm:zeta-suppression-two-point} and in
Hypothesis~\ref{hyp:off-line-toy-signature}.
\end{remark}

% =============================================================================
% Adjacent edits (apply at merge time)
%
% Edit A: §1.2 thm:master-conditional input list.
%   Replace
%     ``the two-point or four-point actual-zeta transverse suppression
%       upper bound (Hypothesis~\ref{hyp:zeta-suppression-two-point} or
%       Hypothesis~\ref{hyp:zeta-suppression-four-point})''
%   with
%     ``the two-point or four-point actual-zeta transverse suppression
%       upper bound
%       (Theorem~\ref{thm:zeta-suppression-two-point} or
%       Hypothesis~\ref{hyp:zeta-suppression-four-point})''
%   and add, after the rigidity-theorem item:
%     ``the off-line toy signature input
%       (Hypothesis~\ref{hyp:off-line-toy-signature})''.
%
% Edit B: §7.1 forward reference cleanup.
%   The previous Hypothesis label hyp:zeta-suppression-two-point is
%   replaced by Theorem label thm:zeta-suppression-two-point.  Search
%   rh_rebuild.tex for occurrences of
%     \ref{hyp:zeta-suppression-two-point}
%   and replace with
%     \ref{thm:zeta-suppression-two-point}
%   (4 occurrences expected: §1.2 master conditional, §6 transverse
%   projection scope remarks, §7 four-point preamble, §8 standalone
%   route remark).
%
% Edit C: §6 transverse projection sufficient-input lemma
%   (lem:transverse-projection-sufficient).
%   Update the closing sentence
%     ``The actual-zeta two-point upper bound must be stated in the
%       same scalar channel''
%   to a present-tense reference:
%     ``The actual-zeta two-point upper bound is stated in the same
%       scalar channel by Theorem~\ref{thm:zeta-suppression-two-point}.''
%
% No other adjacent edits are required.
% =============================================================================
